\documentclass[times,twocolumn,final]{elsarticle}

\usepackage{cag}
\usepackage{framed,multirow}
\usepackage{amssymb}
\usepackage{amsmath}
\usepackage{lineno}
\usepackage{graphicx}
\usepackage{booktabs}
\usepackage{placeins}
\setlength{\emergencystretch}{2em}
\setlength{\hfuzz}{2pt}
\setcounter{totalnumber}{6}
\setcounter{topnumber}{4}
\setcounter{bottomnumber}{4}
\renewcommand{\topfraction}{0.92}
\renewcommand{\textfraction}{0.05}
\renewcommand{\floatpagefraction}{0.85}

\newtheorem{remark}{Remark}

\journal{Computers \& Graphics}

\begin{document}

\begin{frontmatter}

\title{Timestep-Conditioned Adapter Scaling for Multi-view Diffusion Texture Generation}

% Double-anonymized review: author names, affiliations, and acknowledgements
% (including funding) are provided only on the separately uploaded title page.

\begin{abstract}
\iffalse
Multi-view diffusion pipelines for 3D asset texturing commonly inject geometric conditions, such as normal and depth maps, through adapter modules, and the adapter scale at inference time is usually chosen empirically, although the generated views are later baked into the final mesh texture. We show that a single uniform scale is a blunt control: on the studied adapter, aggressive uniform scaling improves structural metrics such as foreground structural similarity (FG-SSIM) but suppresses the texture synthesis capability of the base model, causing flattened surfaces, reduced color variation, and loss of high-frequency details. We further identify an adapter-specific temporal conflict: geometric residuals are most useful for mid-stage structure refinement, whereas strong residual injection in the early or late stages disturbs global formation or suppresses texture synthesis. We introduce a measurable stage-utility analysis with structural guardrails and derive Timestep-Conditioned Adapter Scaling (TCAS), a training-free low--high--low schedule that concentrates strong geometric correction in the middle denoising stage and uses conservative scales otherwise. The schedule C3 $=(1.25, 2.50, 1.25)$ is selected using only a 24-object probe set from 17 scaling variants and then evaluated unchanged, without further search, on the strictly disjoint 276-object holdout. On the full 300-object pooled evaluation with the main adapter, TCAS achieves FG-SSIM comparable to aggressive uniform scaling while improving the peak signal-to-noise ratio (PSNR) by 0.96~dB; it also retains substantially more texture and receives higher preference in a blinded human study. These results suggest treating adapter strength as a stage-dependent inference control variable within the tested MVPainter-style architecture family rather than as a fixed hyperparameter.
\fi
Multi-view diffusion pipelines for 3D asset texturing commonly inject geometric conditions through adapter modules, but adapter scale is not a universal guidance strength. We study Timestep-Conditioned Adapter Scaling (TCAS) as a training-free, three-stage inference intervention: a low--high--low schedule is selected on a probe set and then evaluated under frozen protocols. On a strict 276-object main-adapter holdout, a clean-v2 four-condition run and a separately labeled stage-placement follow-up show that stage placement changes the output, but no schedule dominates all reported targets. In the follow-up, LHL improves on selected targets relative to fixed mean and HLL, whereas LLH is better than LHL on all six non-SSIM evaluation metrics and on the saved-artifact Full-SSIM branch. A second-backbone audit likewise leaves the pre-specified CAI selection undefined. A 12-object stratified baking case study verifies the end-to-end GLB/unseen-view path but does not establish a population-level 3D advantage. We therefore present TCAS as a bounded, adapter-dependent mechanism study rather than a universal schedule-selection rule.
\end{abstract}

\begin{keyword}
\KWD
multi-view diffusion \sep 3D texturing \sep adapter scaling \sep timestep control \sep shape-texture trade-off
\end{keyword}

\end{frontmatter}

\linenumbers

\section{Introduction}
\label{sec:intro}

Recent advances in single-image 3D generation have substantially improved the recovery of geometry, multi-view consistency, and sparse-view reconstruction \cite{liu2023zero1to3,shi2023zero123pp,liu2023syncdreamer,long2023wonder3d,shi2024mvdream,wang2023imagedream,xu2024instantmesh,ke2024era3d}. As geometric generation models become increasingly mature, generating high-quality textures on an existing 3D shape has become a key factor affecting the realism and usability of 3D content. For an object whose geometry has already been determined, the final visual quality depends not only on whether the shape is complete, but also on whether the texture is faithful to the reference image, aligned with the 3D surface, and rich in local details.

Existing multi-view texture generation methods usually generate RGB images from multiple target views using 2D image diffusion models, and then map the generated images back to the 3D surface through projection or texture baking \cite{chen2023text2tex,richardson2023texture,poole2023dreamfusion,chen2023fantasia3d}. This route benefits from the generative capability of image diffusion models, but it still faces challenges in reference appearance alignment, geometry-texture consistency, and local texture detail preservation. Geometry-controlled multi-view diffusion methods represented by MVPainter introduce geometric conditions such as normals and depth maps to improve the alignment between generated textures and the underlying 3D surface, demonstrating the importance of explicit geometric control for 3D texture generation \cite{shao2025mvpainter}.

However, introducing geometric conditions is not only a question of whether a control signal should be used, but also a question of how strongly it should be applied. In adapter-augmented diffusion pipelines, geometric conditions are usually injected into the model through an adapter or control branch, and their influence is regulated by a scaling factor. In practice, this factor is often set empirically: a small scale may lead to insufficient geometric constraints, whereas a larger scale is usually expected to improve structural consistency.

In multi-view texture generation, however, stronger adapter scaling is not necessarily better. Texture generation requires not only contours and structures to be consistent with geometry, but also surface color variations, material details, and high-frequency patterns to be preserved. Overly strong geometric control may force the model to follow normal, depth, or edge conditions too aggressively, thereby suppressing the intrinsic texture synthesis capability of the base diffusion model. In this case, some structural metrics may still improve, while the visual result suffers from texture flattening, reduced color variation, and detail loss.

Therefore, adapter scaling in geometry-controlled multi-view diffusion deserves further analysis. Compared with prior work that mainly studies how to build texture generation pipelines, how to introduce geometric conditions, and how to improve overall generation quality, adapter strength is a fine-grained inference-time control problem. It directly affects the balance between shape consistency and texture detail, and is an easily overlooked factor in multi-view texture generation quality.

This problem is related to, but not solved by, existing guidance scheduling strategies. CFG scheduling changes the extrapolation between conditional and unconditional predictions, while a geometry adapter injects residuals into intermediate representations. Generic ControlNet or adapter scale scheduling similarly treats condition strength mainly as a condition-adherence knob. In texture generation, however, the key question is not only how strongly to follow the geometry, but also when geometric residuals should dominate and when they should recede so that the base model can synthesize color, material, and high-frequency texture. Simple monotonic schedules such as linear or cosine decay do not explicitly address this non-monotonic requirement. Even the recent finding that classifier-free guidance is most useful only in a limited noise interval \cite{kynkanniemi2024limited} does not answer this question: it identifies where conditioning helps for CFG, but not which stages of a geometry-conditioned texture pipeline require strong geometric residuals, where the late-stage risk is texture suppression rather than diversity loss.

The goal of this paper is not to propose a new full texture generation system, but to analyze and improve inference-time adapter scaling within a fixed geometry-controlled multi-view diffusion pipeline.

The main contributions of this paper are as follows:

(1) We systematically diagnose the shape-texture trade-off caused by adapter scaling in multi-view diffusion texture generation, showing that higher scales can improve structural metrics while simultaneously causing adapter-dependent texture degradation.

(2) We formulate Timestep-Conditioned Adapter Scaling (TCAS), a low--high--low inference schedule that concentrates strong geometric correction in the middle denoising stage while using conservative scales otherwise, and test its stage-utility hypothesis with controlled placement comparisons.

(3) We provide a strict-276 main-adapter audit with explicit pre-save versus saved-artifact metric provenance, a fixed-GT Full-SSIM serialization decomposition, and a 17/16/17-step budget audit. The resulting evidence identifies stage sensitivity but also a clear LLH counterexample to a universal LHL optimum.

(4) We assess practical 3D relevance with a traceable 12-object, four-condition GLB baking and unseen-view case study, and audit transfer to a second backbone while keeping cross-backbone scores and undefined CAI outcomes separate.

\section{Related Work}
\label{sec:related}

\subsection{3D Texture Generation and Multi-view Diffusion}
\label{subsec:3dtex}

Single-image 3D generation is often decomposed into geometry generation and texture generation. In recent years, single-image-to-3D, multi-view diffusion, and sparse-view reconstruction methods have made notable progress in structural stability, shape completeness, and cross-category generalization \cite{liu2023zero1to3,shi2023zero123pp,liu2023syncdreamer,long2023wonder3d,shi2024mvdream,wang2023imagedream,xu2024instantmesh,ke2024era3d}. Nevertheless, high-quality texture generation remains a key bottleneck for the realism and usability of 3D content, as reflected by recent surveys of neural 3D mesh texturing \cite{perla2026survey}. Existing texture generation methods usually rely on 2D diffusion models to synthesize RGB images from multiple views, and then project or bake these images onto the 3D surface \cite{chen2023text2tex,richardson2023texture,poole2023dreamfusion,chen2023fantasia3d}. While this strategy inherits the strong generative capability of image diffusion models, it still suffers from reference appearance misalignment, geometry-texture inconsistency, and insufficient local texture detail. Beyond these diffusion pipelines, the graphics community continues to advance surface appearance synthesis more broadly, from exemplar-based texture synthesis \cite{guzman2026adaptive} to diffusion-based procedural material generation \cite{lv2025interpretable}.

To address these issues, methods such as Paint3D, MVPaint, Hunyuan3D 2.0, and MVPainter formulate 3D texture generation as a geometry-controlled multi-view diffusion task, generating multi-view images under geometric control and baking them into UV textures \cite{zeng2024paint3d,cheng2024mvpaint,lai2025hunyuan3d,shao2025mvpainter}. MV-Adapter further shows that adapter-style multi-view control can enhance cross-view consistency without fully retraining the base model \cite{huang2024mvadapter}, and recent texturing systems strengthen consistency during texture synthesis itself, for example through synchronized multi-view diffusion \cite{liu2024synchronized} or consistent-and-seamless text-guided multi-view generation \cite{xu2026wondertex}. By injecting geometric information such as normals and depth maps into the diffusion model as conditional signals, the generated results better follow the 3D surface structure and remain consistent across views. Unlike prior methods that focus mainly on multi-view image generation or general 3D reconstruction, these works emphasize precise alignment between texture and geometry as well as local detail recovery, significantly improving 3D texture generation quality. However, they mainly answer how to introduce geometric conditions and how to construct a texture generation pipeline, while the influence of geometric control strength at inference time on shape and texture quality remains underexplored.

\subsection{Adapter-augmented Diffusion Models and Geometric Control}
\label{subsec:adapter}

Adapter modules have become an important mechanism for introducing additional conditions into diffusion models. T2I-Adapter and ControlNet inject spatial conditions such as edges, depth maps, and normals into a pretrained diffusion model through lightweight parallel branches with a controllable scaling factor \cite{mou2023t2iadapter,zhang2023controlnet}. IP-Adapter introduces image prompts through decoupled attention \cite{ye2023ipadapter}. LoRA and its variants implement lightweight model adaptation through low-rank residual updates \cite{hu2022lora,liu2024dora,zhang2023adalora}. Although these methods differ in architecture, they generally include some form of scaling factor that controls the influence of the external condition or adapter residual on the base model.

In practice, the adapter scale is often treated as an empirical hyperparameter. A larger scale usually improves condition adherence, for example by better following normals, depth maps, or reference images. In texture generation tasks, however, overly strong conditional constraints may weaken the base diffusion model's ability to synthesize color, material, and high-frequency details. Adapter scaling is therefore not a simple matter of increasing the strength as much as possible; it involves a balance between geometric constraints and texture fidelity. This paper focuses on exactly this inference-time control problem: given a fixed adapter structure and fixed geometric conditions, how should the adapter strength be set or scheduled to preserve both shape consistency and texture detail?

\subsection{Diffusion Guidance Scheduling and Timestep Control}
\label{subsec:guidance}

The denoising process of diffusion models is highly stage-dependent. Early steps mainly determine global layout and coarse structure, middle steps refine shape and semantic structure, and late steps are more responsible for color, texture, and high-frequency details \cite{ho2020ddpm,song2021ddim}. Based on this observation, many guidance scheduling strategies have been explored in both research and practice. For example, classifier-free guidance scales may be adjusted linearly, with cosine schedules, or dynamically to mitigate oversaturation, detail loss, or reduced diversity caused by overly strong guidance \cite{ho2022cfg,castillo2023adaptive}. In addition, the development of latent diffusion, high-resolution diffusion, and Transformer-based diffusion models further shows the close relationship among generation quality, noise schedules, conditional injection, and network architecture \cite{rombach2022ldm,esser2024scaling,blackforestlabs2024flux,loshchilov2019decoupled,cherti2023reproducible,nichol2022glide,saharia2022photorealistic,peebles2023scalable}.

These methods all exploit the staged nature of denoising to regulate conditional strength. Notably, Kynk\"a\"anniemi et al.\ show that classifier-free guidance is harmful at high noise levels, largely unnecessary at low noise levels, and beneficial mainly in a middle interval, motivating limited-interval guidance \cite{kynkanniemi2024limited}; mid-stage concentration of conditioning strength therefore has a close precedent in CFG scheduling. However, different control signals act through different mechanisms. CFG scheduling adjusts the weight between conditional and unconditional predictions, mainly affecting the balance among semantic control, diversity, and visual fidelity. By contrast, the scale of a geometric adapter acts directly on intermediate features or residual branches, controlling how strongly geometric conditions intervene in the generation process. TCAS differs from limited-interval guidance in both mechanism and target: it scales adapter feature residuals rather than score extrapolation, raises the middle scale above a conservative baseline instead of switching guidance on and off, and addresses the shape--texture trade-off of geometry-conditioned texture generation, where late-stage over-control suppresses the synthesis of color and high-frequency texture. ControlNet-style scale control is closer in spirit, but it is still usually interpreted as a condition-adherence strength; directly applying a larger, smaller, or smoothly decayed control scale does not specify which denoising stage should receive strong geometric residuals for texture generation.

This distinction matters because the desired schedule for geometry-controlled texture generation is not necessarily monotonic. A monotonic decay or warm-up can reduce over-control in one part of the trajectory, but it cannot simultaneously avoid unstable early residuals, strengthen the middle stage where geometry refinement is most useful, and weaken late residuals that compete with color and high-frequency texture formation. Thus, timestep-conditioned scheduling of adapter strength cannot be treated as a direct transfer of CFG scheduling or generic ControlNet scale scheduling; it requires a dedicated analysis of how geometric residual injection interacts with texture synthesis.

\section{Method}
\label{sec:method}

\subsection{Task Definition and Base Pipeline}
\label{subsec:task}

This paper builds on the MVPainter-style geometry-controlled multi-view texture generation framework \cite{shao2025mvpainter} and studies how the inference-time strength of geometric condition injection affects the generated results. Figure~\ref{fig:overview} illustrates the overall pipeline: the adapter injects geometric residuals into UNet features, TCAS adjusts their strength across denoising stages, and the final multi-view outputs are projected and baked into a textured 3D mesh. Given a single reference image and the corresponding 3D geometry, the base pipeline first generates RGB texture images from six target views and then obtains a textured mesh through projection or texture baking. We do not modify the base UNet, VAE, or pretrained adapter weights. Instead, the adapter scaling factor is treated as the key inference-time control variable.

\begin{figure*}[t]
\centering
\includegraphics[width=\textwidth,trim=0 24pt 0 0,clip]{fig1.pdf}
\caption{Method overview. GeoTex-Adapter generates geometric residuals and injects them into UNet features. TCAS dynamically adjusts the residual strength according to denoising stages to balance geometric consistency and texture fidelity. FAC is an optional learned extension module. The final multi-view outputs are projected and baked into a textured 3D mesh.}
\label{fig:overview}
\end{figure*}

During multi-view diffusion, the 3D geometry is rendered into normal maps, depth maps, and foreground masks, denoted as:
\begin{equation}
\mathbf{G} = \{G_\text{normal}, G_\text{depth}, G_\text{mask}\}.
\end{equation}

The geometric adapter maps these conditions to a residual correction term and injects it into the intermediate hidden states of the UNet. For the $t$-th denoising step, the adapter injection can be written as:
\begin{equation}
\mathbf{h}'_t = \mathbf{h}_t + s(p_t) \cdot A_\phi(\mathbf{h}_t, \mathbf{G}),
\label{eq:injection}
\end{equation}
where $\mathbf{h}_t$ and $\mathbf{h}'_t$ denote the hidden states before and after injection, $p_t$ denotes the sampling progress, and $s(p_t)$ is the adapter scaling strength at the current denoising stage. This formulation highlights a key difference from CFG-style guidance: changing $s(p_t)$ does not reweight two predicted scores, but changes the magnitude of a geometry-conditioned residual added to the feature stream. A larger $s$ usually enhances geometric condition adherence. However, if strong residuals are continuously applied in the late denoising stage, they may suppress the synthesis of color, material, and high-frequency details by the base diffusion model. This paper therefore studies how to schedule $s$ during denoising to balance shape consistency and texture detail preservation.

\subsection{Shape-Texture Diagnostic Metrics}
\label{subsec:metrics}

The influence of adapter scaling on generated results cannot be evaluated by a single structural metric. We divide evaluation metrics into three groups. The first group consists of shape and structure metrics, including FG-SSIM, Edge-SSIM, and Crop-SSIM, which measure consistency with ground-truth renderings in the foreground, edge regions, and local foreground crops based on the structural similarity index \cite{wang2004ssim}. The second group contains signal fidelity and perceptual similarity metrics, including PSNR and the learned perceptual image patch similarity (LPIPS) \cite{zhang2018lpips}, which reflect pixel-level error and differences in deep feature space. The third group consists of texture diagnostic metrics, including Laplacian variance, RGB standard deviation, and gradient magnitude, which are used to detect weakened color variation, surface smoothing, and high-frequency detail loss.

Taking Laplacian variance as an example, the foreground texture response is defined as:
\begin{equation}
M_\text{tex} = \text{Var}(\text{Lap}(I_\text{out} \odot G_\text{mask})),
\end{equation}
where Lap denotes the Laplacian filter, $G_\text{mask}$ denotes the foreground mask, and Var denotes spatial variance. A larger $M_\text{tex}$ usually corresponds to richer local high-frequency variation. If structural metrics improve while Laplacian variance, RGB standard deviation, or gradient magnitude decreases, the improvement may mainly come from contour smoothing or boundary convergence rather than true texture quality improvement.

\subsection{Timestep-Conditioned Adapter Scaling}
\label{subsec:tcas}

The diffusion denoising process is stage-dependent: early steps mainly establish the global layout and coarse contours; middle steps refine geometric structure and cross-view alignment; late steps are more responsible for color, material, and local texture synthesis. Based on this observation, we propose Timestep-Conditioned Adapter Scaling (TCAS), which replaces a globally fixed adapter scale with a stage-dependent scale function. TCAS is not intended as a generic smooth guidance schedule. Its low-high-low form encodes the adapter-specific hypothesis that geometry residuals should be strongest only after coarse structure becomes stable and before late texture synthesis dominates.

To avoid ambiguity caused by different raw timestep directions in different schedulers, we use sampling progress $p$ to represent the denoising process, where $p = 0$ denotes the start from noise and $p = 1$ denotes the end of sampling close to the final image state. TCAS divides denoising into early, middle, and late stages:
\begin{equation}
s(p) = \begin{cases} s_e, & 0 \leq p < 1/3 \\ s_m, & 1/3 \leq p < 2/3 \\ s_l, & 2/3 \leq p \leq 1 \end{cases}.
\label{eq:tcas}
\end{equation}

The selected C3 configuration used in this paper is:
\begin{equation}
(s_e, s_m, s_l) = (1.25, 2.50, 1.25).
\end{equation}

This design follows three principles. First, the early stage uses a conservative scale to avoid disrupting global structure formation with strong geometric residuals. Second, the middle stage uses a higher scale to strengthen geometric constraints from normals, depth maps, and foreground contours. Third, the late stage reduces the scale again to prevent geometric residuals from suppressing color variation, material detail, and high-frequency texture synthesis. This non-monotonic schedule is therefore different from simply decaying or increasing the adapter strength across all timesteps. TCAS introduces no new parameters and does not modify the base model or adapter weights; it only replaces the scaling factor during inference, so its computational overhead is negligible.

\paragraph{How the stage schedule is selected: the CAI interpretation}
We use Correction--Alignment--Interference (CAI) as an empirical diagnostic, not as a theorem or a guaranteed selector. For the main adapter, a 24-object probe was used to nominate C3 before the strict-276 evaluation; the later stage-placement run is explicitly a targeted follow-up. The stage utility notation $U_k=Q(\text{high scale only in }k)-Q(L)$ is useful for describing these results, but additive stage effects are only an approximation because denoising stages interact. Texture diagnostics are treated as variation indicators; fidelity is assessed separately with masked perceptual and color-error metrics. In the independent MV-Adapter audit, the pre-specified CAI rule is set-valued/undefined, so no unique schedule is claimed.

\begin{remark}
The CAI interpretation is adapter-dependent. Changing the adapter, residual magnitude, scale semantics, or backbone requires re-estimating the stage utilities and repeating the probe selection. Confidence intervals describe uncertainty in the measured utilities; they do not establish equivalence or non-inferiority. The independent-backbone experiment therefore tests whether stage-position effects are observable under another implementation, not whether C3 or LHL is a universal choice.
\end{remark}

\subsection{GeoTex-Adapter and FAC Adaptive Extension}
\label{subsec:fac}

GeoTex-Adapter, TCAS, and FAC correspond to three complementary dimensions of adapter control. GeoTex-Adapter determines what geometric residual correction should be injected. TCAS determines when the residual should be strengthened or weakened along the denoising-time dimension. FAC further modulates the residual in a fine-grained manner along spatial positions and frequency components.

GeoTex-Adapter is used as a fixed geometric residual backbone. It receives a 5-channel geometric input consisting of a 3-channel normal map, a 1-channel depth map, and a 1-channel foreground mask. A multi-scale geometric encoder produces a feature pyramid, and residual corrections are injected into multiple resolution levels of the UNet. TCAS acts on the global temporal scale of these residuals to control the influence of geometric conditions at different denoising stages without changing the adapter architecture.

To further examine whether stage-wise adapter control can be extended to learned adaptive correction, we also conduct experiments with Full Adaptive Correction (FAC). FAC contains three lightweight modules: LTAG learns timestep-dependent gating, GSG generates spatial gates from geometric features, and FSC performs frequency-selective correction on adapter residuals. Unlike fixed TCAS, FAC finely modulates adapter residuals in the temporal, spatial, and frequency domains. In the extension experiments, FAC uses TCAS as the base temporal schedule and further adds GSG and FSC for spatial and frequency correction.

It should be emphasized that FAC requires learning a small number of additional parameters. Therefore, TCAS, which requires no training, remains the main method in this paper. FAC is used only as a learned adaptive extension; Section~\ref{subsec:facexp} reports a controlled re-examination that delimits when learned modulation can be expected to help.

\section{Experiments}
\label{sec:experiments}

\subsection{Experimental Setup}
\label{subsec:setup}

\subsubsection{Experimental protocol and comparisons}

Experiments are conducted using an MVPainter-style geometry-controlled multi-view diffusion pipeline. Each object contains a single reference image and corresponding 3D geometry. During preprocessing, the 3D geometry is rendered into target-view normal maps, depth maps, foreground masks, and ground-truth RGB images. The corrected round-two protocol uses six unique target views, $(0,15,12,16,13,14)$ in the forward order (and $(14,15,0,16,12,13)$ after reversal); the legacy order with a duplicated top view is retained only as an audit condition. All comparisons with ground-truth renderings use the same views, $256\times256$ resolution, and foreground masks.

Except for the FAC extension experiment, all scaling experiments keep the object list, reference image, target views, random seed, base UNet, VAE, adapter weights, and number of denoising steps fixed; only the adapter scale or schedule is changed. The number of sampling steps is 50. The historical 24-object probe nominated C3 before the larger evaluation records were generated; the active round-two claims use the strict-276 clean-v2 panel and the separately labeled 276-object stage-placement follow-up. The probe-to-holdout protocol isolates the effect of adapter scaling from changes in model architecture, training data, sampling budget, or evaluation-set tuning.

Accordingly, the comparisons in this paper are comparisons among inference-time adapter scaling strategies under the same base pipeline, rather than a system-level benchmark among different texture generation architectures. Uniform adapter scales serve as the primary deployment baselines because they correspond to the most direct way of choosing a single fixed adapter strength at inference time. The probe variants further test whether the trade-off can be resolved by simple alternatives such as globally high scales, layer-wise redistribution, or moving high adapter strength to early, late, or multiple denoising stages. TCAS changes only how this strength is scheduled across denoising stages.

\subsubsection{Evaluation sets and metrics}

The source protocol records four evaluation sets: (1) a 50-object scale-sweep set; (2) a 26-object texture-audit set; (3) a 24-object probe set for the original schedule nomination; and (4) a 300-object evaluation pool containing a 276-object held-out subset. These exploratory records are retained for provenance; the active round-two results are the strict-276 clean-v2 panel, the separately labeled stage-placement follow-up, and the explicitly bounded extension audits.

Evaluation metrics are divided into three categories: shape and structure metrics (FG-SSIM, Edge-SSIM, and Crop-SSIM), signal and perceptual metrics (PSNR, masked LPIPS, and, where available, CIEDE2000), and texture-variation diagnostics (Laplacian variance, RGB standard deviation, and gradient magnitude). The three variation diagnostics are not direct measures of texture fidelity; their symmetric log-ratio errors relative to ground truth are reported separately. All image metrics are first computed over six unique target views and then averaged over views and objects. For texture-related conclusions, we report means, paired object-level bootstrap intervals, and win rates; DISTS is unavailable in the current baking environment and is not claimed.

The FAC extension uses the same object list and evaluation protocol as the 300-object evaluation pool, and its lightweight modules are trained only on the 1,706-object training pool for the strong-residual adapter instance, which is disjoint from all evaluation objects. TCAS is used as the main baseline, while LTAG, LTAG+GSG, and Full FAC are evaluated as ablation variants under the same paired protocol. The base model and original adapter remain frozen, and only lightweight adaptive correction parameters are learned.

\subsubsection{Implementation and dataset details}

All objects are drawn from the Objaverse 3D asset collection: ground-truth multi-view images are rendered offline with Blender (Cycles engine, 17 views at $256\times256$, white background), and the corrected target-view protocol uses the six unique camera IDs $(0,15,12,16,13,14)$; the legacy duplicate-top order is retained only as an audit condition. The main adapter's training pool contains 1{,}118 rendered objects, while the separately documented strong-residual/FAC audit uses a 1{,}706-object pool from the same source. Both training pools are disjoint from the evaluation pool; we verified that each training/evaluation identifier intersection is zero. Within the 300-object evaluation pool, the 24-object probe set corresponds to the first 24 identifiers ($\mathrm{obj}_{0000}$--$\mathrm{obj}_{0023}$) and the remaining 276 identifiers ($\mathrm{obj}_{0024}$--$\mathrm{obj}_{0299}$) form a holdout that is strictly disjoint from the probe. The exploratory 50-object and 26-object records are retained as provenance; the active evidence freeze uses the strict 276-object cohort and its paired follow-up.

The base pipeline is an MVPainter-style geometry-controlled multi-view diffusion model built on a joint six-view latent UNet, loaded from the local \texttt{MVPainter\_Pipeline} snapshot under \texttt{checkpoints/hf\_repo} with diffusers 0.20.0, and used frozen together with its VAE. The pipeline is reference-image-conditioned; no text prompts are used. The GeoTex-Adapter takes a five-channel geometric input (3-channel normal map, 1-channel depth map, 1-channel foreground mask), produces a multi-scale residual pyramid, and injects corrections into multiple UNet resolution levels; for the adapter studied in the main tables, the requested scale is applied multiplicatively at every injected layer. Inference uses the Euler discrete scheduler with 50 denoising steps and a fixed random seed (42) shared by all schedules within an experiment, so all variants share the same initial latents; the model generates six views per object.

FG-SSIM, Edge-SSIM (edge masks derived from depth/normal discontinuities), PSNR, and LPIPS (AlexNet backbone) are computed on foreground-masked regions after background normalization; the three foreground variation diagnostics are computed on foreground pixels and are not interpreted as texture fidelity without their GT-relative errors. The main tables use GeoTex-Adapter \texttt{geotex\_refattn\_v1/geotex\_step\_0002000.pt} trained on the 1{,}118-object pool. The separate strong-residual/FAC records are labeled as an independent extension protocol and are not pooled with the clean-v2 main-adapter tables. Checkpoint and object-list hashes are recorded in the round-two evidence freeze and reproducibility package.

\paragraph{Statistical reporting}
Object-level paired comparisons report 95\% percentile bootstrap confidence intervals with 10{,}000 resamples, a fixed seed, and object win rates; the direction is reversed for lower-is-better metrics. A confidence interval crossing zero is described as no detected difference, not as evidence of equivalence or non-inferiority. Where a mean is accompanied by $\pm$, the quantity is the standard deviation across objects. The stage-placement follow-up is labeled as such and is not treated as a blind confirmation.

% Historical exploratory sweep/probe tables and qualitative figures are
% retained below only as source provenance. They are not part of the active
% round-two evidence because their checkpoint/metric path predates clean-v2.
\iffalse
\subsection{Shape-Texture Trade-off of Adapter Scaling}
\label{subsec:tradeoff}

Table~\ref{tab:sweep} reports the results of a 13-scale uniform adapter sweep on 50 objects.

\begin{table}[t]
\centering
\caption{Uniform adapter scale sweep results.}
\label{tab:sweep}
\small
\begin{tabular}{lcccc}
\toprule
Scale & $\Delta$FG-SSIM$\uparrow$ & POS & $\Delta$Crop-SSIM$\uparrow$ & $\Delta$PSNR$\uparrow$ \\
\midrule
0.75 & $-$0.013 & 20/50 & +0.113 & 7.53 \\
0.90 & +0.024 & 30/50 & +0.137 & 8.32 \\
1.00 & +0.047 & 34/50 & +0.147 & 8.79 \\
1.10 & +0.070 & 36/50 & +0.159 & 9.08 \\
1.15 & +0.071 & 36/50 & +0.157 & 9.25 \\
1.25 & +0.088 & 37/50 & +0.161 & 9.37 \\
1.35 & +0.100 & 39/50 & +0.163 & 9.63 \\
1.50 & +0.113 & 40/50 & +0.163 & 9.76 \\
1.60 & +0.120 & 42/50 & +0.166 & 9.83 \\
1.75 & +0.127 & 43/50 & +0.173 & 9.97 \\
2.00 & +0.134 & 44/50 & +0.169 & 9.96 \\
2.25 & +0.142 & 43/50 & +0.175 & 9.91 \\
2.50 & +0.145 & 44/50 & +0.165 & 9.91 \\
\bottomrule
\end{tabular}
\end{table}

As the scale increases from $s = 0.75$ to $s = 2.50$, $\Delta$FG-SSIM rises from $-$0.013 to +0.145, and the number of objects with positive improvement increases from 20/50 to 44/50. This indicates that stronger adapter residuals indeed improve the model's adherence to geometric conditions such as normals, depth maps, and foreground contours. Figure~\ref{fig:curve} visualizes this trend: FG-SSIM keeps increasing with adapter scale, while Crop-SSIM saturates and slightly decreases at high scales.

\begin{figure}[t]
\centering
\includegraphics[width=\columnwidth]{fig2.pdf}
\caption{Structural metric changes under uniform adapter scaling. The curves show that FG-SSIM keeps increasing with adapter scale, while Crop-SSIM saturates and slightly decreases at high scales. Stronger geometric control therefore mainly improves structural alignment, but does not necessarily keep improving local foreground quality.}
\label{fig:curve}
\end{figure}

Overall, the FG-SSIM improvement is stable, but the marginal benefit becomes smaller in the high-scale region. For example, increasing the scale from $s = 2.25$ to $s = 2.50$ yields only a limited FG-SSIM gain, while Crop-SSIM already decreases. This suggests that the nature of the high-scale benefit may change: the gain comes more from contour smoothing and foreground alignment than from preservation of local texture details. Figure~\ref{fig:heatmap} shows the per-object $\Delta$FG-SSIM heatmap, confirming that most objects gain higher structural scores as the scale increases, but with clear object-level differences.

\begin{figure}[t]
\centering
\includegraphics[width=\columnwidth]{fig3.pdf}
\caption{$\Delta$FG-SSIM heatmap for 50 objects under 13 adapter scales. Most objects obtain higher structural gains as the scale increases, but there are clear object-level differences.}
\label{fig:heatmap}
\end{figure}

\subsection{Texture Audit: High Scale Causes Texture Flattening}
\label{subsec:audit}

To determine whether the FG-SSIM gain at high scales reflects better visual quality, we further conduct a texture audit on 26 objects, comparing texture metrics and foreground shape changes for $s = 2.50$ relative to $s = 1.25$. Figure~\ref{fig:qualitative} should be read as a local texture comparison rather than as a claim that any method fully matches the ground truth: the key evidence is the progressive loss of color variation, material marks, and high-frequency patterns inside the marked regions as the scale increases. The conservative scale $s = 1.25$ better preserves color variation and local texture, while $s = 2.50$ more easily produces uniform and smooth surfaces.

\begin{figure}[!t]
\centering
\includegraphics[width=\columnwidth]{fig4.pdf}
\caption{Texture degradation under different adapter scales. Each row shows one object; large crops highlight local texture regions (thumbnails below indicate crop locations). Green borders denote conservative scaling $s = 1.25$; red borders denote aggressive scaling $s = 2.50$. Within each row, increasing the adapter scale causes progressive loss of color variation and high-frequency detail: surfaces become more uniform and color-saturated, indicating that stronger geometric correction suppresses the base model's texture synthesis capability.}
\label{fig:qualitative}
\end{figure}

The existing 26-object audit shows a pattern of ``variation loss without a demonstrated shape gain'' for many objects under high scale. Here, a shape gain is not equated with an increase in FG-SSIM: it requires substantive improvement in effective foreground area, contour structure, or geometric detail compared with the conservative scale. If the metric improvement mainly comes from smoothed foreground boundaries or averaged contour errors without visible geometric improvement, it is not categorized as a shape gain. The corrected round-two analysis replaces this audit wording with GT-relative fidelity metrics and held-out evidence.

\begin{figure}[t]
\centering
\includegraphics[width=\columnwidth]{fig5.pdf}
\caption{Texture flattening at $s = 2.50$. Top: RGB standard deviation ratio on 26 objects ($s = 2.50$ / $s = 1.25$). Bottom: gradient magnitude ratio. Values below 1.0 indicate texture loss. Even the eight best-improvement objects (green) show 18\% to 30\% texture loss despite positive FG-SSIM gains.}
\label{fig:audit}
\end{figure}

For multi-view texture generation, overly strong geometric control can suppress the color and detail synthesis capability of the base diffusion model during late denoising steps, making surfaces more uniform. Therefore, scale selection should not rely only on standard metrics such as FG-SSIM or PSNR, but should also incorporate texture-sensitive diagnostics.

\subsection{Timestep-Conditioned Adapter Scaling: Shape-Texture Balance of C3}
\label{subsec:c3}

Based on the texture audit, we do not pursue higher adapter scales throughout the entire denoising process. Instead, we further explore how different scaling strategies balance shape gain and texture loss. We construct 17 adapter scaling variants and evaluate them on the 24-object probe set. To avoid turning the 300-object evaluation pool into a hyperparameter search, the schedule selected in this probe stage is frozen before large-scale evaluation. These variants include three categories: uniform scaling baselines, layer-wise scaling strategies, and timestep-conditioned scaling strategies. Uniform scaling tests the effect of globally strengthening adapter influence and serves as the closest counterpart to using a single ControlNet-style condition scale. Layer-wise scaling examines whether redistributing the same geometric intervention across UNet levels is sufficient. Timestep scaling tests stage placement directly, including variants that put high adapter strength in early, middle, late, or broader temporal regions. This design is intended to rule out the simple explanation that any guidance-like schedule or any stronger control scale would solve the problem.

All variants are evaluated using $\Delta$FG-SSIM for shape gain and $\Delta$Lap Var for texture detail change. Delta texture metrics are computed relative to the corresponding ground-truth foreground renderings under the same views and masks. Importantly, the objective of this probe experiment is not to find the single highest structural score, but to identify a stable compromise strategy that avoids the trap of ``higher structural metrics but flattened texture.'' This criterion is crucial: if a schedule improves FG-SSIM by applying strong geometry residuals late but causes a clear drop in texture response, then it does not solve the adapter-specific conflict even though it appears favorable under a structure-only metric.

\begin{table}[t]
\centering
\caption{Shape-texture trade-off on the 24-object probe set. Delta texture metrics are computed relative to the corresponding ground-truth foreground renderings under the same views and masks; $\Delta$Lap Var closer to 0 indicates smaller texture loss. Eight representative variants of the 17 evaluated are shown; the complete set is visualized in Fig.~\ref{fig:pareto}.}
\label{tab:probe}
\footnotesize\setlength{\tabcolsep}{3pt}
\begin{tabular}{llccc}
\toprule
Variant & Category & $\Delta$FG-SSIM$\uparrow$ & $\Delta$Lap Var $\to$ 0 & Assessment \\
\midrule
C4 & Late-stage high & 0.105 & $-$0.0120 & Trap \\
C2 & Timestep sched. & 0.104 & $-$0.0110 & Trap \\
B2 & Layer-wise sched. & 0.090 & $-$0.0135 & Trap \\
A\_s2.00 & Uniform & 0.087 & $-$0.0062 & Trap \\
C3 (TCAS) & Mid-stage high & 0.084 & $-$0.0006 & Promising \\
A\_s2.25 & Uniform & 0.079 & $-$0.0096 & Trap \\
A\_s2.50 & Uniform & 0.078 & $-$0.0108 & Trap \\
B6 & Layer-wise sched. & 0.000 & $-$0.0027 & OK \\
\bottomrule
\end{tabular}
\end{table}

Table~\ref{tab:probe} summarizes representative scheduling variants on the 24-object probe set. C4 and other late-stage high-scale strategies obtain higher $\Delta$FG-SSIM, but they are accompanied by obvious decreases in Laplacian variance, indicating that their structural gains come at the cost of texture loss. This is the main reason a direct transfer of guidance scheduling intuition is insufficient: a schedule can make the condition stronger at apparently useful timesteps and still damage the texture signal that matters for 3D appearance. Uniform high-scale strategies show a similar problem. For example, $s = 2.50$ achieves $\Delta$FG-SSIM of +0.078, but $\Delta$Lap Var decreases to $-$0.0108. In contrast, C3 achieves $\Delta$FG-SSIM of +0.084. Although it is slightly lower than some aggressive variants in structure gain, its Laplacian variance loss is only $-$0.0006, close to zero texture loss. Therefore, C3 is the only variant labeled Promising. We freeze this configuration before large-scale evaluation and do not tune it on the 300-object set.

The specific configuration of C3 uses the conservative scale $s = 1.25$ in the early and late stages, and the strong scale $s = 2.50$ in the middle geometry-refinement stage, i.e., C3 = (1.25, 2.50, 1.25). Its motivation follows the staged role of diffusion denoising. In the early stage, latent variables are still highly noisy, and strong geometric correction can interfere with global structure formation. In the middle stage, the signal becomes more stable, making it suitable to strengthen geometric constraints such as normals, depth maps, and foreground contours. In the late stage, generation mainly focuses on color, material, and fine texture synthesis, so adapter intervention should be reduced to preserve the texture generation space of the base diffusion model.

Figure~\ref{fig:c3compare} provides a qualitative comparison between C3 and uniform adapter scaling. The red-boxed zoom-in regions make the intended comparison visible: conservative scaling preserves more color variation but provides weaker structural correction, whereas aggressive scaling strengthens geometric constraints but flattens surface texture. C3 applies strong geometric correction in the middle stage and reduces adapter intervention in the late stage, thereby preserving more local texture than $s = 2.50$ while retaining stronger geometric correction than the conservative scale.

\begin{figure}[!t]
\centering
\includegraphics[width=\columnwidth]{fig6.pdf}
\caption{Qualitative comparison between C3 and uniform adapter scaling. Large crops show local texture regions (thumbnails below indicate crop locations). Columns: GT (gray), $s = 1.25$ (green), $s = 2.50$ (red), and C3 (orange). Compared with $s = 2.50$, C3 retains more local color variation and material texture while preserving stronger geometric correction than the conservative scale.}
\label{fig:c3compare}
\end{figure}

Figure~\ref{fig:pareto} further shows the distribution of the 17 variants in the shape-texture space. Most variants with high shape gain are accompanied by large texture loss and fall into the Trap region, whereas C3 maintains a relatively high shape gain while keeping texture loss near zero. Thus, the value of C3 is not that it maximizes every metric, but that it avoids the trap of ``higher structural scores but flattened texture,'' resulting in a more stable shape-texture compromise.

\begin{figure}[!t]
\centering
\includegraphics[width=\columnwidth]{fig7.pdf}
\caption{Shape-texture Pareto trade-off among 17 variants. The horizontal axis is shape gain, where farther right is better; the vertical axis is the Laplacian-variance change, where higher (closer to zero) is better; therefore, the top-right region represents a more desirable compromise. C3 (star) lies in the Promising region, with high shape gain and nearly zero texture loss. Most high-shape-gain variants fall into the Trap region, where structural scores are obtained at the cost of large texture loss. Variants such as B1 have smaller texture degradation but almost no shape gain.}
\label{fig:pareto}
\end{figure}

The results are consistent with a \emph{functional-alignment hypothesis}: geometric residuals may be more useful after coarse structure has emerged and less useful while global layout or fine texture is being formed. This is an interpretation, not an identified causal mechanism. The present 17/16/17 intervention changes both stage placement and nominal residual budget, and the LLH counterexample shows that the data do not identify a unique low--high--low optimum. We therefore treat the stage roles as hypotheses for future budget-matched and content-adaptive tests, not as established causal explanations.

This explanation is consistent with the scheduling-variant results. Early high-scale or early-and-late high-scale variants provide limited structural gain, indicating that early strong intervention is unstable. The late-stage high-scale configuration C4 obtains a high structural score but shows a substantial decrease in Laplacian variance, making it a typical Trap variant where structural metrics improve while texture is flattened. Uniform high-scale variants similarly improve condition adherence but lose texture detail, showing that the issue is not solved by a stronger ControlNet-like scale. An aggressive C3 variant further increases the middle-stage strength but reduces signal fidelity, suggesting that even middle-stage correction has diminishing returns and can lead to over-constraint.

\FloatBarrier

\begin{table}[h!]
\centering
\caption{Comparison of representative variants under extended texture metrics. All delta texture metrics are computed relative to the corresponding ground-truth foreground renderings under the same views and masks. This table verifies that the advantage of C3 is not due to a single Laplacian variance metric.}
\label{tab:extended}
\small
\begin{tabular}{lcccc}
\toprule
Variant & $\Delta$Lap Var & $\Delta$RGB Std & $\Delta$Grad & $\Delta$HF Energy \\
\midrule
C3 & $-$0.001 & $-$0.023 & +0.142 & +0.058 \\
C4 & $-$0.012 & $-$0.032 & +0.086 & +0.062 \\
A\_s2.50 & $-$0.011 & $-$0.044 & +0.054 & +0.078 \\
A\_s2.00 & $-$0.006 & $-$0.040 & +0.100 & +0.070 \\
A\_s1.25 & $-$0.005 & $-$0.032 & +0.104 & +0.053 \\
\bottomrule
\end{tabular}
\end{table}

Table~\ref{tab:extended} compares representative variants under extended texture metrics. C3 has nearly zero loss in Laplacian variance and remains relatively stable in RGB standard deviation, gradient magnitude, and high-frequency energy. Compared with C4 and uniform high-scale strategies, C3 shows weaker texture degradation, indicating that its shape-texture balance is not an artifact of a single metric.

\FloatBarrier

% The historical pooled/CLIP-IQA/preference block below is retained in the
% source archive for provenance only. It is excluded from the revised paper
% because its protocol and saved-artifact provenance do not match the frozen
% clean-v2 evidence. The active results begin with the strict-276 panels.
\fi
\subsection{Strict holdout and stage-placement evidence}
\label{subsec:largescale}

The current main-adapter evidence is reported in two explicitly separated
panels. The first is the strict-276 slice of the clean-v2 four-condition run
(no adapter, fixed-low, fixed-high, and C3), evaluated with the pre-save float
metric path. The second is a targeted stage-placement follow-up on the same
276 object IDs, with fixed-mean, HLL, LHL, and LLH generated in one run. The
repeated C3 outputs are not identical across these two runners (for example,
the mean absolute Full-SSIM difference on the common objects is 0.00395), so
absolute values are not pooled across panels. All paired comparisons within a
panel use the same object IDs and seed.

\begin{table}[t]
\centering
\caption{Strict-276 main-adapter clean-v2 panel. Metrics are pre-save float-evaluation means over the same six unique views; the Full-SSIM column is retained as an audit value and is not silently substituted for the saved-artifact branch.}
\label{tab:strict276}
\scriptsize
\resizebox{\columnwidth}{!}{%
\begin{tabular}{lrrrrrrr}
\toprule
Method & Full-PSNR$\uparrow$ & Full-SSIM$\uparrow$ & Full-LPIPS$\downarrow$ & FG-PSNR$\uparrow$ & FG-SSIM$\uparrow$ & FG-LPIPS$\downarrow$ & Edge-SSIM$\uparrow$ \\
\midrule
No adapter & 10.514 & 0.726 & 0.627 & 8.776 & 0.478 & 0.207 & 0.479 \\
Fixed-low & 15.086 & 0.850 & 0.194 & 7.030 & 0.355 & 0.202 & 0.500 \\
Fixed-high & 13.365 & 0.823 & 0.190 & 5.410 & 0.229 & 0.210 & 0.494 \\
C3 (LHL) & 14.875 & 0.855 & 0.190 & 6.763 & 0.348 & 0.201 & 0.500 \\
\bottomrule
\end{tabular}}
\end{table}

The clean-v2 panel is a non-dominance result: C3 is above fixed-high on
foreground PSNR/SSIM and is close to fixed-low on foreground LPIPS and
Edge-SSIM, but fixed-low has higher Full-PSNR, FG-PSNR, and FG-SSIM. Thus the
main run does not support a blanket claim that TCAS improves practical image
quality over the conservative fixed-low baseline.

\begin{table}[t]
\centering
\caption{Stage-placement follow-up on the strict 276-object cohort. These are the six non-SSIM metrics from one paired run; LHL, HLL, and LLH use 17/16/17 early/middle/late steps. Lower is better for LPIPS.}
\label{tab:stage276}
\scriptsize
\resizebox{\columnwidth}{!}{%
\begin{tabular}{lrrrrrr}
\toprule
Schedule & Full-PSNR$\uparrow$ & Full-LPIPS$\downarrow$ & FG-PSNR$\uparrow$ & FG-SSIM$\uparrow$ & FG-LPIPS$\downarrow$ & Edge-SSIM$\uparrow$ \\
\midrule
Fixed mean & 14.311 & 0.188 & 6.257 & 0.316 & 0.205 & 0.499 \\
HLL & 13.397 & 0.202 & 5.583 & 0.259 & 0.212 & 0.464 \\
LHL (C3) & 14.868 & 0.190 & 6.760 & 0.348 & 0.201 & 0.500 \\
LLH & 15.136 & 0.181 & 6.937 & 0.357 & 0.200 & 0.536 \\
\bottomrule
\end{tabular}}
\end{table}

The follow-up provides evidence that stage placement matters, but it does not
establish LHL as a universal choice. LLH is better than LHL on all six metrics in
Table~\ref{tab:stage276} under their intended directions. This is the central
counterexample to interpreting CAI as a unique LHL selector. The nominal
budgets are also not strictly equal: fixed mean has scale mean 1.6667 and
squared-scale sum 138.8889; HLL and LLH have means 1.675 and squared-scale
sum 157.8125; LHL has mean 1.650 and squared-scale sum 153.125. The recorded
actual residual L2/RMS summaries show further stage-dependent differences, so
the result is a stage-placement sensitivity study rather than a perfectly
equal-effective-budget intervention.

\begin{table}[t]
\centering
\caption{Saved-artifact Full-SSIM in the stage-placement follow-up. Prediction is PNG-reloaded float32; GT is the original RGBA composited over white in float32. C3 deltas use object-level paired bootstrap with 10,000 resamples and seed 20260928.}
\label{tab:serialized276}
\small
\begin{tabular}{lcc}
\toprule
Schedule & Full-SSIM$\uparrow$ & C3 paired difference \\
\midrule
Fixed mean & 0.87395 & +0.00749 \\
HLL & 0.85993 & +0.02151 \\
LHL (C3) & 0.88144 & reference \\
LLH & 0.88313 & $-$0.00170 \\
\bottomrule
\end{tabular}
\end{table}

The corresponding 95\% CIs for C3 minus fixed mean, HLL, and LLH are
$[+0.00674,+0.00826]$, $[+0.01984,+0.02319]$, and
$[-0.00205,-0.00134]$, respectively. The last interval is a detected loss for
C3 on this saved-artifact Full-SSIM branch, not evidence that the schedules
are equivalent. A 12-object saved-tensor decomposition fixes the GT branch
before comparing representations: A fp16 to A fp32 changes mean SSIM by
$+0.03198$; prediction PNG quantization changes it by $-0.00036$; and GT PNG
serialization changes it by $-0.00007$. These diagnostic effects are not used
to retrofit the historical 300-object table.

\FloatBarrier

\subsection{Baking and unseen-view case study}
\label{subsec:baking}

To test whether image-level differences survive into a real 3D output, we
performed a CPU bake on the frozen, stratified 12-object Exact-GLB cohort.
Each of the four main-adapter conditions produced a textured GLB, yielding 48
object--condition exports and 11 unseen-camera renders per export (528
view-level rows). The six source views were fixed and the unseen-view metrics
were aggregated within object before any descriptive comparison.

\begin{table}[t]
\centering
\caption{Descriptive unseen-view baking metrics for the 12-object stratified case study. These values are not a population estimate; DISTS was unavailable.}
\label{tab:bake12}
\small
\resizebox{\columnwidth}{!}{%
\begin{tabular}{lrrrr}
\toprule
Method & Masked PSNR$\uparrow$ & CIEDE2000$\downarrow$ & FG-LPIPS$\downarrow$ & Silhouette IoU \\
\midrule
No adapter & 11.824 & 18.200 & 0.156 & 0.99968 \\
Fixed-low & 6.001 & 48.623 & 0.199 & 0.99968 \\
Fixed-high & 3.891 & 60.420 & 0.216 & 0.99968 \\
C3 (LHL) & 5.638 & 51.183 & 0.201 & 0.99968 \\
\bottomrule
\end{tabular}
}
\end{table}

This case study does not show a C3 advantage over fixed-low or no adapter for
unseen-view fidelity: no adapter has the best descriptive PSNR, CIEDE2000, and
LPIPS means in this cohort. The bake audit verifies semantic mesh/polygon/UV
preservation and embedded textures, but it is not a byte-identity guarantee;
Blender reindexed vertices for \texttt{obj\_0048}. The zero seam values are
reported with the audit's seam definition and denominator and are not treated
as proof of universal seam-free output. Coverage and post-inpainting coverage
are reported separately in the supplementary audit. The run has no generated
12-object GT bake and no DISTS result, and visual inspection found
cross-view/source-fusion colour inconsistency. Thus the 3D result is a
feasibility and limitation case study, not evidence of population-level TCAS
superiority. Supplementary S4 gives object-level failure boundaries:
\texttt{obj\_0078} and \texttt{obj\_0082} have raw coverage below 0.04, while
\texttt{obj\_0048} has 0.8441 inpainted fraction and is the exporter
reindexing case.

\iffalse

Finally, as a transfer test rather than another search, we compare conservative uniform scaling ($s = 1.25$), aggressive uniform scaling ($s = 2.50$), and selected C3 (TCAS) on 300 objects. C3 is not searched again on the 300-object set; it directly uses the timestep schedule selected on the 24-object probe set to evaluate the transferability of the strategy to a larger object set. Because the 24-object probe is a subset of the 300-object pool, the main-table comparison is not probe-disjoint; a disjoint holdout re-verification is provided in Section~\ref{subsec:adapterdep}.

\begin{table}[t]
\centering
\caption{Large-scale pooled evaluation on the full 300-object evaluation pool (pooled statistics; the pool contains the 24 probe objects). Values are absolute metrics of generated outputs with respect to ground truth.}
\label{tab:300obj}
\small\setlength{\tabcolsep}{4pt}
\begin{tabular}{lcccc}
\toprule
Method & FG-SSIM$\uparrow$ & Edge-SSIM$\uparrow$ & PSNR$\uparrow$ & FG-LPIPS$\downarrow$ \\
\midrule
$s = 1.25$ & 0.430 & 0.513 & 17.95 & 0.174 \\
$s = 2.50$ & 0.473 & 0.559 & 17.90 & 0.172 \\
C3 (TCAS) & 0.476 & 0.529 & 18.86 & 0.171 \\
\bottomrule
\end{tabular}
\end{table}


As shown in Table~\ref{tab:300obj}, conservative uniform scaling $s = 1.25$ has lower FG-SSIM and Edge-SSIM, indicating insufficient geometric constraints in the existing audit. Aggressive uniform scaling $s = 2.50$ obtains the highest Edge-SSIM, while C3 has lower Edge-SSIM in the current table; this trade-off is reported explicitly. The numerical comparison in this working draft is not used as the corrected round-two conclusion until the unique-view main-adapter rerun is complete.

PSNR mainly reflects global pixel error and should not be directly equated with texture detail quality. To further verify texture preservation, we compute Laplacian variance, RGB standard deviation, and gradient magnitude on the same 300 objects to measure foreground high-frequency response, color variation amplitude, and local gradient strength, respectively. All results use the same object list, random seed, 50 denoising steps, and adapter weights to ensure comparability.


\begin{table}[t]
\centering
\caption{Large-scale pooled texture-preservation evaluation on the full 300-object evaluation pool (pooled statistics; the pool contains the 24 probe objects). Higher texture metrics indicate better preservation of foreground color variation and high-frequency details. Absolute values are means over objects and views; the ratios in Table~\ref{tab:retention} are the mean over objects of each per-object ratio relative to the corresponding $s = 1.25$ statistic (equal object weighting).}
\label{tab:texture300}
\small
\begin{tabular}{lccc}
\toprule
Method & Lap Var$\uparrow$ & RGB Std$\uparrow$ & Grad Mag$\uparrow$ \\
\midrule
$s = 1.25$ & 0.0151 & 0.0408 & 0.3258 \\
$s = 2.50$ & 0.0095 & 0.0251 & 0.2616 \\
C3 (TCAS) & 0.0129 & 0.0377 & 0.2919 \\
\bottomrule
\end{tabular}
\end{table}

Table~\ref{tab:texture300} reports variation diagnostics only: although $s = 2.50$ achieves the highest Edge-SSIM, it has lower Laplacian variance, RGB standard deviation, and gradient magnitude in this audit. These values indicate a change in foreground variation, not texture fidelity by themselves. The round-two analysis therefore adds symmetric log-ratio errors against ground truth together with masked LPIPS/DISTS and CIEDE2000; the corrected paired holdout values replace this audit-only interpretation when the recovered main-adapter assets are rerun.


\begin{table}[t]
\centering
\caption{Texture retention ratios and object-level stability of C3. Each ratio is the mean over objects of the per-object ratio of the corresponding foreground statistic to that of $s = 1.25$ (equal object weighting); object-level win rates for C3 over $s = 2.50$ are reported in the text.}
\label{tab:retention}
\small\setlength{\tabcolsep}{4pt}
\begin{tabular}{lccc}
\toprule
Method & Lap Var Ratio$\uparrow$ & RGB Std Ratio$\uparrow$ & Grad Mag Ratio$\uparrow$ \\
\midrule
$s = 1.25$ & 1.000 & 1.000 & 1.000 \\
$s = 2.50$ & 0.678 & 0.643 & 0.807 \\
C3 (TCAS) & 0.835 & 0.970 & 0.883 \\
\bottomrule
\end{tabular}
\end{table}


As shown in Table~\ref{tab:retention}, compared with conservative scaling $s = 1.25$, aggressive scaling $s = 2.50$ retains only 67.8\% of Laplacian variance, 64.3\% of RGB standard deviation, and 80.7\% of gradient magnitude, indicating that full high-scale scaling substantially weakens foreground color variation and high-frequency details. In contrast, C3 retains 83.5\%, 97.0\%, and 88.3\% of these metrics, respectively, and clearly outperforms $s = 2.50$ across all three texture measures. This shows that TCAS still incurs some texture reduction, but it substantially alleviates texture flattening compared with full high-scale scaling.

Object-level statistics further show differences between C3 and $s = 2.50$ in the three variation diagnostics and PSNR. These diagnostics are retained as variation/diagnostic metrics rather than labeled texture fidelity. In particular, the revised analysis does not treat a confidence interval crossing zero as proof of equivalence or non-inferiority, and it explicitly reports the Edge-SSIM decrease of C3 relative to $s = 2.50$ together with the corrected paired fidelity metrics.

In summary, the 300-object experiments demonstrate that adapter scaling affects both structural alignment and texture detail preservation. Full high-scale scaling strengthens edge constraints but is more likely to cause texture flattening. TCAS strengthens geometric correction in the middle stage and reduces adapter intervention in the late stage, thereby maintaining structural alignment while substantially reducing texture loss. It should be noted that C3 still has some texture decrease relative to $s = 1.25$; therefore, we do not claim that TCAS completely eliminates texture loss, but rather that it significantly reduces texture degradation compared with full high-scale scaling.

\paragraph{Perceptual evaluation with CLIP-IQA}
To address the concern that proxy texture statistics do not directly measure perceptual quality, we additionally evaluate CLIP-IQA \cite{wang2023clipiqa}, which builds on CLIP image--text representations \cite{radford2021clip}, on a 50-object subset of the evaluation pool as a complementary no-reference perceptual signal. CLIP-IQA is used as complementary perceptual evidence rather than a complete replacement for human preference evaluation.

\begin{table}[t]
\centering
\caption{CLIP-IQA perceptual evaluation on 50 objects from the evaluation pool. Scores are reported for foreground regions and full rendered views; $\pm$ denotes the standard deviation across the 50 objects. Higher is better.}
\label{tab:clipiqa}
\footnotesize\setlength{\tabcolsep}{4pt}
\begin{tabular}{lccc}
\toprule
Method & CLIP-IQA (FG)$\uparrow$ & CLIP-IQA (Full)$\uparrow$ & Win vs. $s=2.50$ \\
\midrule
$s = 1.25$ & $\mathbf{0.4920 \pm 0.0032}$ & $\mathbf{0.4918 \pm 0.0030}$ & $\mathbf{98.0\%}$ \\
$s = 2.50$ & $0.4860 \pm 0.0018$ & $0.4857 \pm 0.0017$ & -- \\
C3 (TCAS) & $0.4906 \pm 0.0028$ & $0.4903 \pm 0.0025$ & 94.0\% \\
\bottomrule
\end{tabular}
\end{table}

Table~\ref{tab:clipiqa} shows that C3 achieves a foreground CLIP-IQA score of 0.4906 versus 0.4860 for $s = 2.50$, scoring higher on 47 of the 50 objects (94.0\%). The paired comparison uses a two-sided paired $t$-test over objects ($t(49)=10.52$, $p\approx3.6\times10^{-14}$), and the 95\% confidence interval of the paired mean difference is constructed by object-level percentile bootstrap with 10{,}000 resamples, giving $[0.0037,\,0.0054]$, which excludes zero. The ordering $s = 1.25$ $>$ C3 $>$ $s = 2.50$ supports the diagnosis that aggressive scaling reduces perceptual quality, while TCAS recovers most of this loss with stronger geometry than conservative scaling.

\paragraph{Blinded preference study}
To further examine whether the diagnostic texture differences are perceptually meaningful, we conducted a blinded three-alternative forced-choice (3AFC) preference study on 30 objects drawn from the evaluation pool. Twenty-four participants viewed, for each object and each of the three criteria (texture naturalness, shape consistency, overall quality), the anonymized outputs of conservative uniform scaling, aggressive uniform scaling, and C3 side by side; the left-to-right order of the three conditions was randomized independently per trial, and participants were not informed which method produced which image or how many conditions existed. Instructions asked participants to select the best image for the stated criterion only. Each participant judged all 30 objects for all three criteria, yielding 720 valid first-choice decisions for each criterion. The study complements the proxy diagnostics by testing whether the observed texture and structure differences are reflected in human preference.

\begin{table}[!htbp]
\centering
\caption{Blinded three-alternative first-choice preference study. Percentages are computed over 720 valid decisions per criterion.}
\label{tab:preference}
\footnotesize\setlength{\tabcolsep}{3.5pt}
\begin{tabular}{lccc}
\toprule
Method & Texture naturalness$\uparrow$ & Shape consistency$\uparrow$ & Overall quality$\uparrow$ \\
\midrule
$s = 1.25$ & $\mathbf{332/720\ (46.1\%)}$ & 96/720 (13.3\%) & 176/720 (24.4\%) \\
$s = 2.50$ & 104/720 (14.4\%) & $\mathbf{331/720\ (46.0\%)}$ & 126/720 (17.5\%) \\
C3 (TCAS) & 284/720 (39.4\%) & 293/720 (40.7\%) & $\mathbf{418/720\ (58.1\%)}$ \\
\bottomrule
\end{tabular}
\end{table}
\FloatBarrier

The preference pattern in Table~\ref{tab:preference} follows the intended trade-off. Conservative scaling is most often selected for texture naturalness, while aggressive scaling is most often selected for shape consistency. C3 stays close to both single-scale baselines on their respective strengths and receives the highest overall preference, with 58.1\% of all overall-quality votes. Thus, TCAS is not claimed to maximize every individual criterion; it provides a perceptually preferred balance between geometric adherence and texture preservation.

Because first-choice decisions are clustered within both participants and objects, we additionally report 95\% confidence intervals obtained by a two-level cluster bootstrap (resampling participants and objects with replacement, 10{,}000 resamples). C3's overall-quality share is 58.1\% (95\% CI $[50.1\%, 65.8\%]$), exceeding the conservative and aggressive baselines by $+33.6$ (CI $[19.3, 47.2]$) and $+40.6$ percentage points (CI $[29.4, 51.8]$); both differences exclude zero, so the overall-preference advantage of C3 is significant under cluster-aware inference. By contrast, the baselines' single-criterion advantages are not significant (texture naturalness: $s = 1.25$ over C3 $+6.7$ points, CI $[-7.9, 21.5]$; shape consistency: $s = 2.50$ over C3 $+5.3$ points, CI $[-6.7, 17.9]$). Overall quality is therefore the only criterion with a significant preference, supporting the interpretation of TCAS as the perceptually preferred balance rather than a per-criterion maximum.

\fi

\subsection{FAC Adaptive Correction Extension: A Controlled Re-examination}
\label{subsec:facexp}

TCAS uses a fixed three-stage scale, which has the advantages of requiring no training, being simple to implement, and introducing low computational overhead. However, its residual scaling is still uniform across spatial positions and frequency components, lacking fine-grained adaptivity. To examine whether stage-wise adapter control can be extended to learned adaptive correction, we trained the FAC extension---LTAG for timestep gating, GSG for spatial gating, and FSC for frequency-selective correction on top of the TCAS temporal schedule. The base model, VAE, and GeoTex-Adapter remain frozen; only the three lightweight controllers are trainable, adding roughly $6\times10^{3}$ parameters. The controllers are trained with the same enhanced denoising objective as the base adapter---a latent noise-prediction MSE with $1.7\times$ foreground and $1.3\times$ edge spatial reweighting plus a latent-SSIM term---using AdamW (peak learning rate $10^{-4}$, 100 warm-up steps) for 2{,}000 steps on the adapter's training pool, which is strictly disjoint from all evaluation objects, and all variants are evaluated under the separately documented strong-residual, per-layer-capped audit protocol with identical objects, seeds, and evaluation code.

In this controlled setting, no learned variant reached the training-free schedule. We tested seven configurations spanning joint, envelope-preserving, fidelity-weighted, and gently scheduled controllers. Their paired deficits relative to the TCAS baseline range from $-1.1$ to $-2.9$~dB foreground PSNR (FG-PSNR; win rates at most 16\%). A warm-start control in which the gate is left untrained at the exact C3 envelope shows no detected difference from TCAS (paired FG-PSNR $\Delta=-0.18$~dB on the 12-object control, 95\% CI $[-0.42,+0.05]$, crossing zero); this interval does not prove equivalence or validate transfer. No configuration or hyperparameter was selected on any evaluation object; the complete dose--response study is provided in the supplementary material.

For transparency, we note that an initial evaluation of the FAC extension, conducted on the main adapter under our earlier protocol before the strictly paired and strictly disjoint training controls were adopted, had suggested a positive gain. That initial comparison differed in both evaluation protocol and adapter instance, and it does not meet the controls we now consider necessary for a transfer claim; under the strictly paired, strictly disjoint protocol above, no trained configuration replicates the gain. We therefore report only the controlled re-examination and treat the initial comparison as superseded.

The outcome delimits the tested extension and does not identify a causal mechanism: every trained configuration degraded the measured fidelity under this protocol, while the untrained envelope control did not show a detected difference. Learning a modulation that improves upon C3 therefore requires a training signal that protects foreground fidelity; this is future work. The FAC results are reported as a negative, separate extension study, not as a claim that TCAS is the only viable method.

\FloatBarrier

\iffalse
\subsection{Adapter-Dependence of High-Scale Failure Modes}
\label{subsec:adapterdep}

The preceding experiments demonstrate the texture-flattening failure mode under high adapter scaling on the studied adapter checkpoint. A natural question is whether this failure mode is universal or specific to the adapter's trained residual properties. We conduct additional analysis across adapter checkpoints with different residual magnitudes to characterize how the failure mode depends on adapter properties.

We find that the specific failure mode of aggressive adapter scaling is not universal. Across three adapter checkpoints spanning a spectrum of trained residual magnitudes, uniform high scale ($s = 2.50$) produced: (a) texture flattening when all layers are scaled without per-layer caps; (b) high-frequency artifact amplification when per-layer caps restrict the fine-texture layer to a low scale while coarse layers are boosted; or (c) no appreciable effect when the adapter's residual magnitude is very small. A per-layer scale-cap ablation on the same checkpoint reproduced both (a) and (b), isolating the inference-time scale semantics as the determining factor for the failure \emph{direction}.

A residual-magnitude sweep further shows that the intervention \emph{strength} is monotonic in the trained residual magnitude: scaling the adapter's output projection from 0.1$\times$ to 3$\times$ transitions the behavior from no-effect, to artifact amplification under high scale, to over-constraint even under conservative scale.

Moreover, a six-object stage ablation shows that even \emph{which denoising stage is most sensitive to high scale} is adapter-dependent: on the artifact-prone adapter, early-stage high scale reproduced nearly all of the damage (FG-SSIM $\downarrow$ from 0.266 to 0.187) while late-stage high scale was nearly harmless (FG-SSIM 0.261), whereas on the paper's main adapter the late stage was the primary source of texture loss.

Despite this variability, in every regime where scaling had a non-negligible effect, the temporal low-high-low schedule C3 yielded the best shape-texture balance. This indicates that TCAS's value is structural rather than incidental: by relying on the diffusion process's stage-dependent roles (global layout $\to$ geometry refinement $\to$ texture synthesis), it mitigates the harm of high-scale intervention across the tested adapter regimes, despite their adapter-specific failure modes. The failure mode (flatten vs.\ artifact vs.\ no-op), its strength, and its temporal location are all adapter-dependent; what remains stable in our experiments is that the middle denoising stage was the best-performing tested window for geometric correction, although whether this holds across adapter families remains an open question for future work.

Two further robustness checks were performed on the strong-residual, per-layer-capped adapter instance described above. First, a boundary/peak sensitivity sweep (three high-scale windows $[0.25,0.75)$, $[1/3,2/3)$, $[0.4,0.6)$ $\times$ five peak scales from 2.00 to 3.00, fifteen candidates on the 24-object probe protocol) shows that the canonical middle-third window lies in the stable high-performance region: its 2.50 setting achieves 16.71~dB whole-image PSNR, the best setting in the same window (2.75) achieves 16.79~dB with worse texture diagnostics and lower FG-SSIM, and the broader and narrower windows reach at most 16.73 and 16.62~dB. C3 is therefore a robust middle-stage choice rather than a boundary-sensitive point estimate. Second, because the 24-object probe is part of the 300-object pool, transfer claims were re-verified on the disjoint 276-object holdout ($\mathrm{obj}_{0024}$--$\mathrm{obj}_{0299}$), summarized in Table~\ref{tab:holdoutsched}: C3 improves PSNR over the strongest non-C3 probe candidates by $+0.258$~dB (95\% CI $[0.207,0.308]$, 206/276 wins) over a trapezoid schedule and $+0.324$~dB ($[0.281,0.367]$, 232/276 wins) over a Gaussian-peak schedule, and over the reviewer-relevant generic schedules by $+0.323$~dB ($[0.263,0.386]$, 214/276 wins, FG-SSIM statistically tied) over linear warm-up and $+0.902$~dB ($[0.824,0.981]$, 253/276 wins, FG-SSIM $+0.046$) over the cosine bump. The holdout ordering matches the probe: warm-up approaches C3 in structure but loses in signal fidelity, and the smooth bump loses in both because its effective high-scale duration is shorter than the full middle third.

\begin{table}[t]
\centering
\caption{Scheduling transfer on the disjoint 276-object holdout ($\mathrm{obj}_{0024}$--$\mathrm{obj}_{0299}$; strong-residual, per-layer-capped adapter instance of Section~\ref{subsec:adapterdep}). $\Delta$PSNR (dB) is the paired C3-minus-competitor difference (positive favors C3); 95\% CIs use object-level percentile bootstrap with 10{,}000 resamples. Absolute metrics are means over holdout objects and views.}
\label{tab:holdoutsched}
\footnotesize\setlength{\tabcolsep}{2.5pt}
\begin{tabular}{lcccc}
\toprule
Schedule & FG-SSIM$\uparrow$ & PSNR$\uparrow$ & $\Delta$PSNR vs.\ C3 (95\% CI) & C3 wins \\
\midrule
Trapezoid & 0.356 & 19.04 & $+0.258$ $[0.207, 0.308]$ & 206/276 \\
Gaussian peak & 0.358 & 18.97 & $+0.324$ $[0.281, 0.367]$ & 232/276 \\
Linear warm-up & 0.377 & 18.97 & $+0.323$ $[0.263, 0.386]$ & 214/276 \\
Cosine bump & 0.330 & 18.39 & $+0.902$ $[0.824, 0.981]$ & 253/276 \\
C3 (TCAS) & 0.376 & 19.29 & -- (reference) & -- \\
\bottomrule
\end{tabular}
\end{table}


\fi

\subsection{Second-backbone audit and contribution boundary}
\label{subsec:adapterdep}

We audited the same stage-position question on an independent MV-Adapter
implementation using an Exact-GLB cohort of 76 objects and 11 recorded
conditions. The full table and paired bootstrap output are provided in the
reproducibility package. Dynamic schedules with high scale 1.50 and 1.00 are
kept as separate scale families; their rows are not pooled. The table below
shows the paired schedule comparison and the conservative controls.

\begin{table}[t]
\centering
\caption{Selected within-backbone MV-Adapter means on the Exact 76-object cohort. Absolute scores are not compared with the MVPainter-style backbone.}
\label{tab:mvadapter}
\scriptsize
\resizebox{\columnwidth}{!}{%
\begin{tabular}{lrrrrrr}
\toprule
Method & PSNR$\uparrow$ & FG-SSIM$\uparrow$ & Edge-SSIM$\uparrow$ & FG-LPIPS$\downarrow$ & $\Delta E_{00}\downarrow$ & GT-texture$\downarrow$ \\
\midrule
No geometry & 13.114 & 0.498 & 0.233 & 0.183 & 20.720 & 1.701 \\
Fixed-low & 13.357 & 0.528 & 0.240 & 0.188 & 20.150 & 2.224 \\
Fixed mean & 13.337 & 0.528 & 0.240 & 0.188 & 20.197 & 2.203 \\
HLL & 13.321 & 0.527 & 0.240 & 0.188 & 20.243 & 2.232 \\
LHL & 13.344 & 0.528 & 0.240 & 0.188 & 20.182 & 2.269 \\
LLH & 13.353 & 0.528 & 0.240 & 0.188 & 20.157 & 2.199 \\
\bottomrule
\end{tabular}}
\end{table}

The audit supports an observable stage-position effect within this backbone,
but not a uniform practical gain. LHL has slightly higher FG-SSIM and
Edge-SSIM than HLL and LLH in the paired schedule comparison, while LLH has the
highest PSNR among those four and lower GT-relative texture error than LHL.
LHL has no clear advantage over fixed-low or fixed mean under the paired
confidence intervals. The frozen CAI rule is
\texttt{undefined\_set\_valued}, and the
official MV-Adapter pretraining UID list was unavailable, so pretraining
disjointness is not asserted. We therefore treat this result as a within-
backbone diagnostic, not as a cross-backbone causal validation or a unique
schedule-selection result.

\section{Limitations}
\label{sec:limitations}

This paper mainly focuses on inference-time adapter scaling and does not retrain the base multi-view diffusion model or the adapter. The evidence supports a bounded stage-placement effect within the tested protocols, not a universal optimizer. The main clean-v2 run does not dominate fixed-low, and the stage follow-up favors LLH over LHL on all six non-SSIM metrics and on saved-artifact Full-SSIM.

Second, Laplacian variance, RGB standard deviation, and gradient magnitude are proxy metrics for texture degradation. They are effective for detecting high-frequency loss and surface smoothing, but they cannot fully replace human perceptual evaluation. For objects with strong reflection, transparency, thin structures, or highly stylized textures, discrepancies may exist between these metrics and subjective quality.

In addition, the current TCAS uses fixed stage boundaries and a fixed scale combination, C3 = (1.25, 2.50, 1.25). The actual 50-step partition is 17/16/17, and the nominal scale sums, squared sums, and residual norms differ across schedules; consequently the placement result is not a strict equal-effective-budget causal test. The second-backbone CAI rule is undefined/set-valued, and its official pretraining UID disjointness is unknown. The schedule has not been validated across seeds, sampler families, object complexity, material type, or reference-image quality. Future work could study budget-matched stage controls, content-adaptive boundaries, and pre-registered scale-selection rules.

Finally, our controlled FAC re-examination is a separate strong-residual extension protocol. It is negative under the paired controls and is not used to claim that TCAS is universally superior to learned modulation.


\section{Conclusion}
\label{sec:conclusion}

This paper presents a controlled analysis of adapter scaling in multi-view diffusion texture generation. The evidence shows that stage placement changes both structural and fidelity metrics, and that the effects are adapter- and protocol-dependent. Stronger control can improve selected structural targets while degrading other fidelity or texture measures.

We formulate Timestep-Conditioned Adapter Scaling (TCAS) as a simple low--high--low inference intervention. The 24-object probe nominated C3 for the main adapter, while the strict-276 clean-v2 and stage-placement panels provide the current evidence boundary. Within the follow-up, LHL improves over fixed mean and HLL on several targets, but LLH is better on all six non-SSIM metrics and on saved-artifact Full-SSIM.

Therefore, the defensible contribution is a reproducible, training-free stage-utility/mechanism study: it exposes a shape--texture trade-off, supplies a concrete residual-scaling implementation, and documents both positive stage effects and counterexamples. It does not establish a unique CAI selector, a universal LHL optimum, cross-backbone causality, or a population-level 3D quality advantage. The results motivate treating adapter strength as a stage-dependent control variable whose useful schedule must be re-estimated for each adapter and protocol.

The 12-object baking study confirms that the full GLB and unseen-view path is operational, but its descriptive metrics favor no adapter over C3 and fixed-low for this stratified cohort. This negative practical result is retained as part of the paper's scope rather than hidden as a selection effect.

\section*{Data availability}

All data included in this study are available upon request by contact with the corresponding author.

\section*{Declaration of competing interests}

The authors declare that they have no known competing financial interests or personal relationships that could have appeared to influence the work reported in this paper.


\begin{thebibliography}{45}
\bibitem{liu2023zero1to3}
Liu R, Wu R, Van Hoorick B, et al. Zero-1-to-3: Zero-shot one image to 3D object. \textit{Proceedings of the IEEE/CVF International Conference on Computer Vision}, 2023.

\bibitem{shi2023zero123pp}
Shi R, Chen H, Zhang Z, et al. Zero123++: A single image to consistent multi-view diffusion base model. arXiv:2310.15110, 2023.

\bibitem{liu2023syncdreamer}
Liu Y, Lin C, Zeng Z, et al. SyncDreamer: Generating multiview-consistent images from a single-view image. \textit{International Conference on Learning Representations}, 2024.

\bibitem{long2023wonder3d}
Long X, Guo Y C, Lin C, et al. Wonder3D: Single image to 3D using cross-domain diffusion. \textit{Proceedings of the IEEE/CVF Conference on Computer Vision and Pattern Recognition}, 2024.

\bibitem{shi2024mvdream}
Shi Y, Wang P, Ye J, et al. MVDream: Multi-view diffusion for 3D generation. \textit{International Conference on Learning Representations}, 2024.

\bibitem{wang2023imagedream}
Wang P, et al. ImageDream: Image-prompt multi-view diffusion for 3D generation. arXiv:2312.02201, 2023.

\bibitem{xu2024instantmesh}
Xu J, et al. InstantMesh: Efficient 3D mesh generation from a single image with sparse-view large reconstruction models. \textit{Advances in Neural Information Processing Systems}, 2024.

\bibitem{ke2024era3d}
Li P, Liu Y, Long X, et al. Era3D: High-resolution multiview diffusion using efficient row-wise attention. \textit{Advances in Neural Information Processing Systems}, 2024.

\bibitem{chen2023text2tex}
Chen D Z, et al. Text2Tex: Text-driven texture synthesis via diffusion models. \textit{Proceedings of the IEEE/CVF International Conference on Computer Vision}, 2023.

\bibitem{richardson2023texture}
Richardson E, Metzer G, Alaluf Y, et al. TEXTure: Text-guided texturing of 3D shapes. \textit{ACM SIGGRAPH 2023 Conference Proceedings}, 2023.

\bibitem{poole2023dreamfusion}
Poole B, Jain A, Barron J T, Mildenhall B. DreamFusion: Text-to-3D using 2D diffusion. \textit{International Conference on Learning Representations}, 2023.

\bibitem{chen2023fantasia3d}
Chen R, Chen Y, Jiao N, Jia K. Fantasia3D: Disentangling geometry and appearance for high-quality text-to-3D content creation. \textit{Proceedings of the IEEE/CVF International Conference on Computer Vision}, 2023.

\bibitem{shao2025mvpainter}
Shao M, Xiong F, Sun Z, Xu M. MVPainter: Accurate and detailed 3D texture generation via multi-view diffusion with geometric control. arXiv:2505.12635, 2025.

\bibitem{kynkanniemi2024limited}
Kynk\"a\"anniemi T, Aittala M, Karras T, et al. Applying guidance in a limited interval improves sample and distribution quality in diffusion models. \textit{Advances in Neural Information Processing Systems}, 2024.

\bibitem{perla2026survey}
Perla S R K, Zhang H, Mahdavi-Amiri A. Advances in neural 3D mesh texturing: A survey. \textit{Computer Graphics Forum}, 2026, e70392.

\bibitem{guzman2026adaptive}
Guzm\'an J E, Mould D, Paquette E. Adaptive multiresolution exemplar-based texture synthesis on animated fluids. \textit{Computers \& Graphics}, 2026, 134: 104490.

\bibitem{lv2025interpretable}
Lv X, Wu Z, Xu J, et al. Interpretable procedural material graph generation via diffusion models from reference images. \textit{The Visual Computer}, 2025, 41(13): 11195--11205.

\bibitem{zeng2024paint3d}
Zeng X, Chen X, Qi Z, et al. Paint3D: Paint anything 3D with lighting-immersive texture generation. \textit{Proceedings of the IEEE/CVF Conference on Computer Vision and Pattern Recognition}, 2024.

\bibitem{cheng2024mvpaint}
Cheng W, Mu J, Zeng X, et al. MVPaint: Synchronized multi-view diffusion for painting anything 3D. \textit{Proceedings of the IEEE/CVF Conference on Computer Vision and Pattern Recognition}, 2025.

\bibitem{lai2025hunyuan3d}
Zhao Z, Lai Z, Lin Q, et al. Hunyuan3D 2.0: Scaling diffusion models for high resolution textured 3D assets generation. arXiv:2501.12202, 2025.

\bibitem{huang2024mvadapter}
Huang Z, et al. MV-Adapter: Multi-view consistent image generation made easy. \textit{Proceedings of the IEEE/CVF International Conference on Computer Vision}, 2025.

\bibitem{liu2024synchronized}
Liu Y X, Xie M S, Liu H Y, Wong T T. Text-guided texturing by synchronized multi-view diffusion. \textit{SIGGRAPH Asia 2024 Conference Papers}, 2024.

\bibitem{xu2026wondertex}
Xu Q, Han X G, Zhang L. WonderTex: Consistent-and-seamless texture generation with text-guided multi-view image diffusion models. \textit{IEEE Transactions on Visualization and Computer Graphics}, 2026, 32(7): 5815--5825.

\bibitem{mou2023t2iadapter}
Mou C, Wang X, Xie L, et al. T2I-Adapter: Learning adapters to dig out more controllable ability for text-to-image diffusion models. \textit{Proceedings of the AAAI Conference on Artificial Intelligence}, 2024.

\bibitem{zhang2023controlnet}
Zhang L, Rao A, Agrawala M. Adding conditional control to text-to-image diffusion models. \textit{Proceedings of the IEEE/CVF International Conference on Computer Vision}, 2023.

\bibitem{ye2023ipadapter}
Ye H, Zhang J, Liu S, Han X, Yang W. IP-Adapter: Text compatible image prompt adapter for text-to-image diffusion models. arXiv:2308.06721, 2023.

\bibitem{hu2022lora}
Hu E J, Shen Y, Wallis P, et al. LoRA: Low-rank adaptation of large language models. \textit{International Conference on Learning Representations}, 2022.

\bibitem{liu2024dora}
Liu S Y, Wang C Y, Yin H, et al. DoRA: Weight-decomposed low-rank adaptation. \textit{International Conference on Machine Learning}, 2024.

\bibitem{zhang2023adalora}
Zhang Q, Zuo S, Liang C, et al. AdaLoRA: Adaptive budget allocation for parameter-efficient fine-tuning. \textit{International Conference on Learning Representations}, 2023.

\bibitem{ho2020ddpm}
Ho J, Jain A, Abbeel P. Denoising diffusion probabilistic models. \textit{Advances in Neural Information Processing Systems}, 2020.

\bibitem{song2021ddim}
Song J, Meng C, Ermon S. Denoising diffusion implicit models. \textit{International Conference on Learning Representations}, 2021.

\bibitem{ho2022cfg}
Ho J, Salimans T. Classifier-free diffusion guidance. \textit{NeurIPS Workshop on Deep Generative Models and Downstream Applications}, 2021. arXiv:2207.12598.

\bibitem{castillo2023adaptive}
Castillo A, Kohler J, P\'erez J C, et al. Adaptive guidance: Training-free acceleration of conditional diffusion models. \textit{Proceedings of the AAAI Conference on Artificial Intelligence}, 2025.

\bibitem{rombach2022ldm}
Rombach R, Blattmann A, Lorenz D, Esser P, Ommer B. High-resolution image synthesis with latent diffusion models. \textit{Proceedings of the IEEE/CVF Conference on Computer Vision and Pattern Recognition}, 2022.

\bibitem{esser2024scaling}
Esser P, Kulal S, Blattmann A, et al. Scaling rectified flow transformers for high-resolution image synthesis. \textit{International Conference on Machine Learning}, 2024.

\bibitem{blackforestlabs2024flux}
Black Forest Labs. FLUX: Flow matching for image generation. https://github.com/black-forest-labs/flux, 2024.

\bibitem{loshchilov2019decoupled}
Loshchilov I, Hutter F. Decoupled weight decay regularization. \textit{International Conference on Learning Representations}, 2019.

\bibitem{cherti2023reproducible}
Cherti M, Beaumont R, Wightman R, et al. Reproducible scaling laws for contrastive language-image learning. \textit{Proceedings of the IEEE/CVF Conference on Computer Vision and Pattern Recognition}, 2023.

\bibitem{nichol2022glide}
Nichol A, Dhariwal P, Ramesh A, et al. GLIDE: Towards photorealistic image generation and editing with text-guided diffusion models. \textit{International Conference on Machine Learning}, 2022.

\bibitem{saharia2022photorealistic}
Saharia C, Chan W, Saxena S, et al. Photorealistic text-to-image diffusion models with deep language understanding. \textit{Advances in Neural Information Processing Systems}, 2022.

\bibitem{peebles2023scalable}
Peebles W, Xie S. Scalable diffusion models with transformers. \textit{Proceedings of the IEEE/CVF International Conference on Computer Vision}, 2023.

\bibitem{wang2004ssim}
Wang Z, Bovik A C, Sheikh H R, Simoncelli E P. Image quality assessment: From error visibility to structural similarity. \textit{IEEE Transactions on Image Processing}, 2004, 13(4): 600--612.

\bibitem{zhang2018lpips}
Zhang R, Isola P, Efros A A, Shechtman E, Wang O. The unreasonable effectiveness of deep features as a perceptual metric. \textit{Proceedings of the IEEE/CVF Conference on Computer Vision and Pattern Recognition}, 2018.

\bibitem{wang2023clipiqa}
Wang J, Chan K C K, Loy C C. Exploring CLIP for assessing the look and feel of images. \textit{Proceedings of the AAAI Conference on Artificial Intelligence}, 2023.

\bibitem{radford2021clip}
Radford A, Kim J W, Hallacy C, et al. Learning transferable visual models from natural language supervision. \textit{International Conference on Machine Learning}, 2021.

\end{thebibliography}
\end{document}
